\chapter*{برای مطالعه‌ی بیشتر}
\addcontentsline{toc}{chapter}{برای مطالعه‌ی بیشتر}
\markboth{برای مطالعه‌ی بیشتر}{برای مطالعه‌ی بیشتر}

منابع زیر برای خواننده‌ای انتخاب شده‌اند که بخواهد بعضی از بحث‌های
کتاب را با جزئیات بیشتری دنبال کند. این فهرست با کتاب‌نامه‌ی اصلی
تفاوت دارد: کتاب‌نامه آثاری را ثبت می‌کند که در متن به آن‌ها ارجاع
داده شده است؛ این بخش مسیرهایی برای مطالعه‌ی عمیق‌تر پیشنهاد می‌کند.

\begin{enumerate}
  \item \fullcite{shannon1948}. متن بنیادین نظریه‌ی ریاضی ارتباط؛ برای خواندن صورت‌بندی اصلیِ خوداطلاعاتی، آنتروپی و ظرفیت کانال.
  \item \fullcite{pierce1980}. مقدمه‌ای تاریخی و مهندسی‌محور بر نمادها، سیگنال‌ها، نویز و مسئله‌ی انتقال اطلاعات.
  \item \fullcite{cover2006}. کتاب درسی جامع برای مطالعه‌ی دقیق‌تر آنتروپی، اطلاعات متقابل، کدگذاری و نظریه‌ی کانال.
  \item \fullcite{li-vitanyi2008}. مرجع تخصصی برای پیچیدگی کولموگروف، توصیف الگوریتمی و نسبت آن با مفهوم اطلاعات.
  \item \fullcite{sep-peirce-semiotics}. مدخلی دانشگاهی برای آشنایی عمیق‌تر با نظریه‌ی نشانه‌ها و سه‌گانه‌ی پیرس.
  \item \fullcite{schrodinger1944}. متنی تاریخی درباره‌ی مسئله‌ی نظم و وراثت در سلول زنده؛ باید آن را در زمینه‌ی علمی زمان انتشارش خواند.
  \item \fullcite{oppenheim2010}. مرجعی دانشگاهی برای تحلیل سیگنال‌های پیوسته و گسسته و پردازش سیگنال در زمان گسسته.
  \item \fullcite{openstax-physics}. منبع آموزشی آزاد برای مرور پتانسیل الکتریکی، جریان، قانون اهم و مفاهیم پایه‌ی فیزیک.
\end{enumerate}
